%=====================================================================
\chapter{Sampling and Study Power}\label{ch:sampling}
%=====================================================================
\locator{4}{Part~4 --- Measurement, Data and Analysis}

\begin{outcomes}
By the end of this chapter you will be able to:
\begin{tight}
  \item \textbf{Relate} the four quantities --- $\alpha$, power, effect size,
    $N$ --- and solve for any one.
  \item \textbf{Choose} a smallest effect size of interest and defend it.
  \item \textbf{Run} an a-priori power analysis and report it.
  \item \textbf{Explain} why post-hoc power is uninformative and why
    underpowered studies inflate published effects.
  \item \textbf{Estimate} sample size for qualitative work without pretending
    it is a power calculation.
\end{tight}
\end{outcomes}

\begin{prereq}
\chref{experiments} for designs, \chref{survey} for sampling frames, and
\chref{measurement} --- power depends on measurement precision, so an unreliable
measure costs you sample size.
\end{prereq}

\begin{keyterms}
type I and type II error $\bullet$ $\alpha$ $\bullet$ power $\bullet$ effect
size $\bullet$ Cohen's $d$ $\bullet$ Cliff's $\delta$ $\bullet$ smallest effect
size of interest $\bullet$ a-priori power analysis $\bullet$ post-hoc power
$\bullet$ winner's curse $\bullet$ design effect $\bullet$ saturation
\end{keyterms}

\section{Four Quantities, Three Free}

\begin{center}
\small
\begin{tabular}{@{}L{2.6cm}L{11.7cm}@{}}
\toprule
$\alpha$        & Probability of a false positive: concluding there is an
                  effect when there is none. Conventionally 0.05, which is a
                  convention, not a law \\
Power $1-\beta$ & Probability of detecting an effect that is really there.
                  Conventionally 0.80, meaning you accept a one-in-five chance
                  of missing a real effect \\
Effect size     & How big the effect is, in standardised or natural units \\
$N$             & Sample size, at the unit of randomisation \\
\bottomrule
\end{tabular}
\end{center}

Fix any three and the fourth is determined. In practice $\alpha$ and power are
set by convention, so the real question is always the trade between effect size
and $N$: \emph{how small an effect do I need to be able to detect, and can I
afford the sample that requires?}

\begin{alertbox}[Do this before the study, not after]
An a-priori power analysis is a design decision. A post-hoc one --- computing
power from the effect you observed --- is uninformative, and specifically so:
it is a deterministic function of the p-value. A non-significant result always
yields low ``observed power'', so reporting it adds nothing beyond restating
that the result was not significant.

\medskip
If you want to say something useful after a null result, report the confidence
interval on the effect, or run an equivalence test (\chref{inference}). Those
answer the question post-hoc power is usually reached for --- ``was my study
capable of detecting anything worth detecting?'' --- and they answer it
correctly.
\end{alertbox}

\section{Choosing an Effect Size, Honestly}

This is where power analysis is usually corrupted. The scholar needs $N$ to
come out affordable, so they choose a large expected effect, and the
calculation obligingly returns a small sample.

\begin{center}
\small
\begin{tabular}{@{}L{4.0cm}L{10.3cm}@{}}
\toprule
\textbf{Basis} & \textbf{Assessment} \\
\midrule
Meta-analytic estimate from the literature
  & Best available. Note that published effects are inflated by publication
    bias, so shade downwards \\
A pilot study
  & \textbf{Dangerous.} Effect estimates from small pilots are extremely noisy,
    and if the pilot happened to over-estimate, you power for an effect that
    does not exist. Pilots are for feasibility and instrumentation, not effect
    size \\
Smallest effect size of interest
  & \textbf{Best practice.} Not what you expect, but the smallest effect that
    would matter to anyone. Below it you do not care whether the effect exists \\
Cohen's conventions (0.2 / 0.5 / 0.8)
  & A last resort. Cohen offered them for fields with no information and
    warned against routine use \\
\bottomrule
\end{tabular}
\end{center}

\begin{conceptbox}[The smallest effect size of interest turns statistics into an argument about substance]
Instead of ``what effect do I expect?'', ask ``how large would the effect have
to be for anyone to act on it?''

\medskip
If a code review technique saved 30 seconds per review, no team would change
their process; a saving of 10 minutes would matter. That judgement makes the
target effect a claim about practice, defensible in the thesis and arguable by
an examiner --- which is exactly what you want. It also protects you: a study
powered for the smallest interesting effect gives a meaningful null result,
because failing to find an effect that size is worth knowing.
\end{conceptbox}

\subsection{Common effect sizes}

\begin{center}
\small
\begin{tabular}{@{}L{3.0cm}L{4.6cm}L{6.7cm}@{}}
\toprule
\textbf{Measure} & \textbf{For} & \textbf{Notes} \\
\midrule
Cohen's $d$
  & Difference between two means
  & $d = (\bar{x}_1 - \bar{x}_2)/s_{\text{pooled}}$. Assumes roughly normal,
    similar variances \\
Hedges' $g$
  & As $d$, small samples
  & Corrects $d$'s upward bias when $N$ is small. Prefer it below about 20 per
    group \\
Cliff's $\delta$
  & Ordinal or non-normal data
  & Probability one group exceeds the other, minus the reverse.
    Distribution-free; the right default for software engineering data, which
    is rarely normal \\
Odds ratio / risk ratio
  & Binary outcomes
  & Report the risk ratio too: odds ratios are routinely misread as risk
    ratios and overstate effects when outcomes are common \\
$r$, $R^2$
  & Association, variance explained
  & $R^2$ is sample-dependent; report adjusted $R^2$ and an interval \\
\bottomrule
\end{tabular}
\end{center}

\begin{worked}[Powering the code review experiment]
\textbf{The design.} Two review techniques compared on defects identified.
Between-subjects, $\alpha = 0.05$, two-tailed, power 0.80.

\medskip
\textbf{Step 1: smallest effect of interest.} Reviews currently find about 6
defects on average, with a standard deviation of roughly 2.5. The team lead
judges that fewer than 1 extra defect per review would not justify changing
the process. So the SESOI is 1 defect, giving
$d = 1 / 2.5 = 0.40$.

\medskip
\textbf{Step 2: required N.}
\begin{lstlisting}[language=Rlang]
library(pwr)
pwr.t.test(d = 0.40, sig.level = 0.05, power = 0.80,
           type = "two.sample", alternative = "two.sided")
#   n = 99.08 per group  ->  100 per group, 200 total
\end{lstlisting}

\medskip
\textbf{Step 3: confront reality.} You can recruit 30 developers. At $n = 15$
per group, power to detect $d = 0.40$ is about 0.14 --- a 14\% chance of
finding a real effect of the size that matters. Running this study is close to
pointless.

\medskip
\textbf{Step 4: the four honest responses.}
\begin{tight}
  \item \textbf{Change the design.} Within-subjects: each developer reviews
    under both techniques. This removes between-subject variance and
    dramatically increases power. With a correlation of about 0.6 between a
    developer's two scores, roughly 40 to 50 developers-worth of paired
    observations suffice --- and 30 developers each doing four reviews may
    reach it, if the nesting is modelled (\chref{regression}).
  \item \textbf{Change the outcome.} A more precise measure --- fewer missed
    defects on a seeded-defect task, where the ground truth is known --- shrinks
    the standard deviation and raises $d$ for the same real effect.
  \item \textbf{Change the question} to an exploratory or qualitative one
    (\chref{qualitative}), which 30 people amply supports.
  \item \textbf{Run it as a registered pilot} whose declared purpose is to
    estimate parameters for a later multi-site study, and report it as such.
\end{tight}

\medskip
\textbf{What is not honest} is running it at $n = 15$, obtaining $p = 0.22$,
and writing ``no significant difference was found between the techniques''.
That sentence implies evidence of no difference, and the study had almost no
capacity to find one.

\medskip
Figures illustrative.
\end{worked}

\section{Why Underpowered Studies Damage the Literature}

An underpowered study is not merely uninformative. It is actively harmful, and
the mechanism is worth understanding~\cite{button2013power}.

\begin{center}
\small
\begin{tabular}{@{}L{3.6cm}L{10.7cm}@{}}
\toprule
\textbf{Consequence} & \textbf{Mechanism} \\
\midrule
Inflated published effects (the winner's curse)
  & With low power, only unusually large sample effects reach significance. So
    the effects that get published are systematically over-estimates ---
    sometimes by a factor of two or more \\
Low positive predictive value
  & Among all significant results in a low-powered literature, a large
    proportion are false positives, even with $\alpha = 0.05$ \\
Failed replications that teach nothing
  & A replication powered for the inflated published effect is itself
    underpowered for the true one \\
Wasted participants
  & Subjects gave time to a study that could not answer its question --- which
    is an ethical issue, not only a methodological one (\chref{ethics}) \\
\bottomrule
\end{tabular}
\end{center}

\begin{conceptbox}[The consolation for a small field]
Software engineering research often cannot recruit hundreds of professionals.
That is a real constraint, and the answer is not to give up on quantitative
work. It is to:

\begin{tight}
  \item prefer within-subject designs, which buy power for free;
  \item use precise measures, since precision substitutes for sample size;
  \item report effect sizes with confidence intervals rather than dichotomous
    significance verdicts, so that a small study contributes an estimate even
    when it cannot settle a question;
  \item and design so the result can be \emph{pooled}. Three well-reported
    small studies with compatible measures can be meta-analysed
    (\chref{slr}); three studies reporting only p-values cannot.
\end{tight}

\medskip
That last point is the most important one in this chapter. A small study that
reports its effect size, its interval and its raw data has contributed
permanently to the evidence base. One that reports only ``no significant
difference'' has contributed nothing.
\end{conceptbox}

\section{Sample Size in Qualitative Work}

Power analysis does not apply, and pretending otherwise is a category error.
But ``as many as I could get'' is not an answer either.

\begin{center}
\small
\begin{tabular}{@{}L{3.6cm}L{10.7cm}@{}}
\toprule
\textbf{Basis} & \textbf{How to state it} \\
\midrule
Saturation
  & With a stopping rule declared in advance and evidence of new codes per
    interview (\chref{qualitative}) \\
Information power
  & Fewer participants are needed when the aim is narrow, the sample is
    specific, the theory is established and the dialogue is strong. State which
    of these apply \\
Design requirement
  & A cross-case design needs enough cases for the comparison: theoretical
    replication requires at least one of each type (\chref{casestudy}) \\
Pragmatic
  & Access limited you. Say so plainly, and discuss what it means for the
    claim. This is an acceptable answer honestly given \\
\bottomrule
\end{tabular}
\end{center}

\begin{pitfall}[Five power mistakes]
\begin{tight}
  \item \textbf{Post-hoc power} reported after a null result.
  \item \textbf{Powering from a pilot's observed effect}, which is noise.
  \item \textbf{Powering at the wrong unit} --- 200 reviews from 20 developers
    is not $N = 200$ (\chref{experiments}).
  \item \textbf{Ignoring the design effect} in clustered sampling: correlated
    units mean the effective sample is smaller than the count.
  \item \textbf{Ignoring measurement unreliability.} An unreliable measure
    attenuates the observed effect, so the $N$ you need is larger than the
    calculation assuming perfect measurement suggests.
\end{tight}
\end{pitfall}

\begin{traditions}{sample size}
\tradEMP{Determined by a-priori power analysis against the smallest effect of
interest, at the unit of randomisation, and reported in full.}
\tradDSR{Evaluation sample size follows the FEDS strategy: a naturalistic
summative evaluation may involve one organisation, and the claim is scoped
accordingly (\chref{designscience}).}
\tradQUAL{Determined by saturation or information power, demonstrated rather
than asserted. More participants is not automatically better; depth competes
with breadth for the same hours.}
\tradFORM{Not applicable. A proof holds for all instances, which is precisely
its advantage.}
\end{traditions}

\begin{lab}[Power your own study, then face the answer]
\begin{tightnum}
  \item State your smallest effect size of interest in \emph{natural units},
    and justify it with a sentence about practice.
  \item Convert to a standardised effect using a variance estimate from the
    literature or your pilot.
  \item Compute required $N$ with \texttt{pwr} for your design.
  \item Compare with what you can actually recruit.
  \item If short, work through the four responses in the worked example and
    choose one. Write down which, and why.
\end{tightnum}

\medskip
Step 5 is the deliverable. Portfolio piece P3 asks for exactly this decision,
and Appendix~\ref{app:schedule} gives credit for abandoning an infeasible
design on evidence.
\end{lab}

\begin{checkpoint}
\begin{tightnum}
  \item Why is post-hoc power uninformative, and what should be reported
    instead after a null result?
  \item Explain the winner's curse and why it makes an underpowered literature
    worse than no literature.
  \item Why is powering from a pilot's observed effect size dangerous?
  \item Your measure has reliability 0.6 rather than 0.9. Which direction does
    that move your required sample size, and why?
\end{tightnum}
\end{checkpoint}

\begin{chaptersummary}
\begin{tight}
  \item $\alpha$, power, effect size and $N$: fix three, the fourth follows.
    The real trade is effect size against $N$.
  \item Do power analysis before, not after. Post-hoc power is a function of
    the p-value and adds nothing.
  \item Power for the smallest effect size of interest, not the expected one.
    That makes the target a defensible claim about practice.
  \item Cliff's $\delta$ is a better default than Cohen's $d$ for software
    engineering data, which is rarely normal.
  \item Underpowered studies inflate published effects and fill the literature
    with false positives.
  \item When you cannot recruit: within-subject designs, precise measures,
    effect sizes with intervals, and reporting that permits later pooling.
  \item Qualitative sample size rests on saturation or information power,
    demonstrated --- or on honestly stated access constraints.
\end{tight}
\end{chaptersummary}

\begin{reviewq}
  \item Power of 0.80 accepts a one-in-five failure to detect a real effect,
    while $\alpha = 0.05$ accepts one-in-twenty false positives. Why are these
    asymmetric, and is the asymmetry defensible for software engineering
    research?
  \item Suppose every software engineering experiment were required to be
    powered at 0.80 for a defensible SESOI. What would happen to the field's
    published output, and would the field be better off?
  \item Design a multi-site study for your topic that would achieve adequate
    power by combining several institutions' small samples, and identify the
    methodological problems that combination introduces.
  \item Underpowered studies waste participants' time, which \chref{ethics}
    treats as an ethical matter. Should ethics review boards assess statistical
    power? Argue both sides.
\end{reviewq}
